\chapter*{Abstract} % senza numerazione
\label{sommario}

\addcontentsline{toc}{chapter}{Abstract} % da aggiungere comunque all'indice

Conversational artificial intelligence on edge devices must satisfy three
constraints at once: a per-token generation latency low enough for natural
spoken dialogue, a single-digit-watt power envelope, and local-only execution
for privacy. Meeting these constraints with a modern transformer-based language
model on hardware with limited resources is non-trivial, because such models require significant compute, memory, and power.


This thesis investigates whether a combination of aggressive quantization and
selective hardware acceleration can bring a competitive conversational model
within the edge envelope. The target model is the 1.2-billion-parameter language
backbone of Liquid AI's LFM2.5-Audio, an end-to-end speech model co-designed for
efficient CPU inference. Fine-grained profiling of the model on a quad-core
Arm Cortex-A53 shows that a single kernel, the MatMul + SwiGLU feed-forward
block, accounts for about 68\% of the per-token runtime. This result motivates a
heterogeneous design in which only that dominant kernel is offloaded to a custom
FPGA accelerator on the AMD-Xilinx Kria KV260, while attention, normalization,
the gated short convolution, and the final projection remain on the CPU.

The accelerator is implemented in Vitis HLS as a five-stage dataflow pipeline
operating on 4-bit (Q4\_K) weights and 8-bit activations, and is integrated into
the \texttt{llama.cpp}/\texttt{ggml} inference stack through a fused hardware
operator with permanent per-layer weight caching, interrupt-driven
synchronization, and transparent fallback to CPU execution. The design is
analysed at stage granularity, showing that performance is bound by off-chip
memory bandwidth rather than arithmetic capacity, and that it operates near the
device's logic and routing limits (87\% LUT occupancy).

The work includes the model profiling, the
software integration in \texttt{llama.cpp}, the FPGA accelerator design and
characterization, a throughput/efficiency study across CPU thread counts
(T1--T4), and a comparison against an STM32MP257F-DK edge MPU. The accelerated
path improves performance-per-watt by up to 57\% (prefill) and 44\% (decode) at
low thread counts and enables sub-500\,ms decode latency with only two active
CPU threads, while the analysis quantifies the memory-bandwidth, orchestration,
and routing constraints that bound the attainable throughput.
